
\subsubsection{Benchmark Fragmentation}

Meta-analysis of 7,635 benchmarks across 4,886 LLM evaluation papers found less than 25\% overlap in benchmark usage across studies. This fragmentation creates evaluation infrastructure where researchers cannot easily compare results across papers. Our work extends this observation by testing whether leaderboard-derivable statistics can prospectively identify reliable benchmarks within this fragmented ecosystem.

\subsubsection{Cross-Benchmark Ranking Disagreement}

Studies have demonstrated systematic ranking disagreement across 37 LLM models evaluated on 6 metric sets, showing that model rankings are sensitive to evaluation methodology. Our hypothesis proposes CV as a predictor of this stability, operationalized as mean cross-benchmark Spearman $\rho$. The null result refutes this relationship but reveals negative correlations between benchmarks, suggesting disagreement may arise from valid construct divergence, not just measurement noise.

\subsubsection{Construct Validity in Psychometrics}

Campbell \& Fiske's (1959) multitrait-multimethod framework distinguishes convergent validity (measures of the same construct should correlate highly) from discriminant validity (measures of different constructs should correlate weakly or negatively). Applied to LLM benchmarks, negative cross-benchmark correlations could indicate good discriminant validity---benchmarks measure distinct trust sub-dimensions---or poor convergent validity---benchmarks claiming to measure ''trust'' don't agree on what that means. This ambiguity is critical: our CV-stability hypothesis assumed low cross-benchmark $\rho$ indicates measurement unreliability (noise), but psychometric theory reveals an alternative explanation---low $\rho$ may reflect valid construct orthogonality.
